\textbf{Phần 1. Trắc nghiệm lựa chọn}
\setcounterpageref{ex}{0}

\begin{ex} %[4H6N28]
Để vẽ đường thẳng đi qua điểm $O$ và vuông góc với đường thẳng $a$, bước đầu tiên ta đặt
\choice
{Thước thẳng bất kì}
{Com-pa tại điểm $O$}
{\True Một cạnh góc vuông của ê-ke trùng với đường thẳng $a$}
{Thước đo góc tại $O$}
\loigiai{
Đặt một cạnh góc vuông của ê-ke trùng với đường thẳng $a$.\\
\textbf{Đáp án của câu hỏi là C}
}
\end{ex}

\begin{ex} %[4H6N28]
Khi vẽ đường cao $AH$ của tam giác $ABC$, ta vẽ đường thẳng đi qua đỉnh $A$ và
\choice
{Song song với cạnh $BC$}
{\True Vuông góc với cạnh $BC$}
{Đi qua trung điểm của $BC$}
{Cắt cạnh $AB$}
\loigiai{
Đường cao kẻ từ $A$ vuông góc với cạnh đáy $BC$.\\
\textbf{Đáp án của câu hỏi là B}
}
\end{ex}

\begin{ex} %[4H6H28]
Hình tam giác có bao nhiêu đường cao?
\choice
{\True $3$ đường cao}
{$1$ đường cao}
{$2$ đường cao}
{$4$ đường cao}
\loigiai{
Mỗi tam giác có 3 đỉnh nên có 3 đường cao.\\
\textbf{Đáp án của câu hỏi là A}
}
\end{ex}

\begin{ex} %[4H6H28]
Trong tam giác vuông $ABC$ (vuông tại $A$), hai đường cao xuất phát từ đỉnh $B$ và đỉnh $C$ chính là
\choice
{Đường trung tuyến}
{Cạnh huyền $BC$}
{Không tồn tại}
{\True Hai cạnh góc vuông $BA$ và $CA$}
\loigiai{
Hai cạnh góc vuông chính là hai đường cao xuất phát từ hai đỉnh nhọn.\\
\textbf{Đáp án của câu hỏi là D}
}
\end{ex}

\begin{ex} %[4H6H28]
Để vẽ hình chữ nhật có chiều dài $5\text{ cm}$ và chiều rộng $3\text{ cm}$, ta cần vẽ các góc tại 4 đỉnh là
\choice
{Góc nhọn}
{Góc tù}
{\True Góc vuông}
{Góc bẹt}
\loigiai{
4 đỉnh của hình chữ nhật đều là góc vuông.\\
\textbf{Đáp án của câu hỏi là C}
}
\end{ex}

\begin{ex} %[4H6H28]
Dùng ê-ke để kiểm tra góc của một khung ảnh hình vuông, ta thấy góc đó
\choice
{Bé hơn góc vuông ê-ke}
{Lớn hơn góc vuông ê-ke}
{Không đo được}
{\True Khớp hoàn toàn với góc vuông của ê-ke}
\loigiai{
Khung ảnh hình vuông có các góc vuông $90^\circ$.\\
\textbf{Đáp án của câu hỏi là D}
}
\end{ex}

\begin{ex} %[4H6V28]
Vẽ đường cao tương ứng với đáy $CD$ của hình thang $ABCD$, ta vẽ đoạn thẳng hạ từ đỉnh $A$ hoặc $B$
\choice
{Song song với $CD$}
{Đi qua trung điểm $CD$}
{\True Vuông góc với đáy $CD$}
{Bằng cạnh bên}
\loigiai{
Đường cao vuông góc với cạnh đáy.\\
\textbf{Đáp án của câu hỏi là C}
}
\end{ex}

\begin{ex} %[4H6V28]
Cho điểm $M$ nằm ngoài đường thẳng $d$. Có thể vẽ được bao nhiêu đường thẳng đi qua $M$ và vuông góc với $d$?
\choice
{$0$}
{\True $1$}
{$2$}
{Vô số}
\loigiai{
Chỉ có duy nhất 1 đường thẳng.\\
\textbf{Đáp án của câu hỏi là B}
}
\end{ex}

\begin{ex} %[4H6V28]
Để vẽ một hình vuông cạnh $4\text{ cm}$, ta cần thực hiện vẽ bao nhiêu góc vuông liên tiếp?
\choice
{\True $4$ góc vuông (hoặc ít nhất 3 góc vuông)}
{$2$ góc vuông}
{$1$ góc vuông}
{$0$ góc vuông}
\loigiai{
Hình vuông có 4 góc vuông.\\
\textbf{Đáp án của câu hỏi là A}
}
\end{ex}

\begin{ex} %[4H6V28]
Khi vẽ đường cao của tam giác có góc tù ở đáy, chân đường cao nằm ở
\choice
{Nằm bên trong cạnh đáy}
{\True Nằm trên phần kéo dài của cạnh đáy}
{Trùng với một đỉnh}
{Không vẽ được}
\loigiai{
Chân đường cao nằm trên đường thẳng chứa cạnh đáy (phần kéo dài).\\
\textbf{Đáp án của câu hỏi là B}
}
\end{ex}

\begin{ex} %[4H6V28]
Để gấp một tờ giấy hình chữ nhật thành một góc vuông, ta có thể
\choice
{\True Gấp đôi mép thẳng của tờ giấy sao cho hai mép trùng nhau}
{Gấp mép giấy bất kì}
{Vò nhàu tờ giấy}
{Cắt xéo tờ giấy}
\loigiai{
Gấp đôi mép thẳng của tờ giấy tạo thành góc bẹt chia đôi là góc vuông.\\
\textbf{Đáp án của câu hỏi là A}
}
\end{ex}

\begin{ex} %[4H6C28]
Cho tam giác $ABC$. Qua mỗi đỉnh $A, B, C$ ta lần lượt vẽ đường thẳng vuông góc với cạnh đối diện. Ba đường thẳng này sẽ
\choice
{Song song với nhau}
{Tạo thành tam giác mới}
{Không bao giờ cắt nhau}
{\True Cùng cắt nhau tại một điểm (trực tâm)}
\loigiai{
Ba đường cao của một tam giác luôn đồng quy tại một điểm (gọi là trực tâm).\\
\textbf{Đáp án của câu hỏi là D}
}
\end{ex}

\textbf{Phần 2. Trắc nghiệm đúng sai}
\setcounterpageref{ex}{0}

\begin{ex} %[4H6H28]
Xét tính đúng sai về thao tác thực hành vẽ:
\choiceTF
{\True Sử dụng ê-ke để vẽ đường thẳng vuông góc với một đường thẳng cho trước.}
{\True Mỗi tam giác có 3 đường cao tương ứng với 3 cạnh đáy.}
{Tam giác vuông không có đường cao nào.}
{\True Trong tam giác vuông, hai cạnh góc vuông chính là hai đường cao.}
\loigiai{
\begin{itemize}
    \item Dùng ê-ke vẽ góc vuông: Đúng.
    \item Tam giác có 3 đường cao: Đúng.
    \item Tam giác vuông vẫn có 3 đường cao. Sai.
    \item Hai cạnh góc vuông là 2 đường cao: Đúng.
\end{itemize}
\begin{itemize}
    \item a) \textbf{Đúng}.
    \item b) \textbf{Đúng}.
    \item c) \textbf{Sai}.
    \item d) \textbf{Đúng}.
\end{itemize}
}
\end{ex}

\begin{ex} %[4H6V28]
Cho đoạn thẳng $AB = 6\text{ cm}$. Điểm $I$ là trung điểm của $AB$. Qua $I$ vẽ đường thẳng $d$ vuông góc với $AB$. Xét tính đúng sai:
\choiceTF
{\True Đoạn thẳng $IA = IB = 3\text{ cm}$.}
{\True Đường thẳng $d$ vuông góc với $AB$ tại điểm $I$.}
{Góc tạo bởi $d$ và $AB$ có số đo là $60^\circ$.}
{\True Đường thẳng $d$ là đường trung trực của đoạn thẳng $AB$.}
\loigiai{
\begin{itemize}
    \item $IA = IB = 3\text{ cm}$: Đúng.
    \item $d \perp AB$ tại $I$: Đúng.
    \item Góc tạo bởi $d$ và $AB$ là $90^\circ \neq 60^\circ$. Sai.
    \item Đường thẳng vuông góc tại trung điểm là đường trung trực: Đúng.
\end{itemize}
\begin{itemize}
    \item a) \textbf{Đúng}.
    \item b) \textbf{Đúng}.
    \item c) \textbf{Sai}.
    \item d) \textbf{Đúng}.
\end{itemize}
}
\end{ex}

\textbf{Phần 3. Trả lời ngắn}
\setcounterpageref{ex}{0}

\begin{ex} %[4H6H28]
Một hình tam giác có bao nhiêu đường cao?
\shortans{3}
\loigiai{
\textbf{Đáp số:} $3$\\
3 đường cao.
}
\end{ex}

\begin{ex} %[4H6H28]
Số góc vuông của thước ê-ke là bao nhiêu?
\shortans{1}
\loigiai{
\textbf{Đáp số:} $1$\\
1 góc vuông.
}
\end{ex}

\begin{ex} %[4H6V28]
Qua một điểm $A$ ở ngoài đường thẳng $d$, vẽ được bao nhiêu đường thẳng vuông góc với $d$?
\shortans{1}
\loigiai{
\textbf{Đáp số:} $1$\\
1 đường thẳng duy nhất.
}
\end{ex}

\begin{ex} %[4H6V28]
Một hình chữ nhật có bao nhiêu góc vuông?
\shortans{4}
\loigiai{
\textbf{Đáp số:} $4$\\
4 góc vuông.
}
\end{ex}

\begin{ex} %[4H6V28]
Một hình thang có thể có nhiều nhất bao nhiêu góc vuông?
\shortans{2}
\loigiai{
\textbf{Đáp số:} $2$\\
2 góc vuông (hình thang vuông).
}
\end{ex}

\begin{ex} %[4H6C28]
Trong một tam giác đều, số đo mỗi góc là bao nhiêu độ?
\shortans{60}
\loigiai{
\textbf{Đáp số:} $60$\\
$180 : 3 = 60^\circ$.
}
\end{ex}
